\chapter{The Explicit Formula and the Riemann Hypothesis}\label{ch:explicit}

\begin{center}\itshape
If I were to awaken after having slept for a thousand years, my first question would be: has the Riemann hypothesis been proven?\\
\upshape --- attributed to David Hilbert
\end{center}

At $x=10$ the prime powers are $2,3,4,5,7,8,9$. Each contributes its $\Lambda$-value, so
\[
\psi(10)=\log2+\log3+\log2+\log5+\log7+\log2+\log3=\log2520=7.832.
\]
The number $10$ is not a prime power, so $\psi_0(10)=\psi(10)$. The function $\psi$ is a staircase: it stays flat between prime powers and jumps by $\log p$ at $p^k$. The explicit formula writes the halfway value of that staircase as a straight line minus one wave for each zero of $\zeta$ in the critical strip.

The identity that is proved, Corollary~\ref{cor:exactexplicit}, is for $\psi_0$, with the sum taken symmetrically:
\[
\psi_0(x)=x-\lim_{T\to\infty}\sum_{|\gamma|<T}\frac{x^\rho}\rho-\log2\pi-\frac12\log\Bigl(1-\frac1{x^2}\Bigr).
\]
At a prime power the halfway convention costs half the jump: $\psi(p^k)-\psi_0(p^k)=\frac12\log p$. The series $\sum 1/|\rho|$ diverges, so the symmetric limit belongs in the statement. A conjugate pair $\rho,\overline\rho$ contributes $-2\operatorname{Re}(x^\rho/\rho)$. With $\rho=\beta+i\gamma$ and $x=e^u$, that contribution is a wave in the variable $u=\log x$, of amplitude $2e^{\beta u}/|\rho|$ and frequency $\gamma$.

Figure~\ref{fig:explicit} draws the staircase against two truncations of the right-hand side. Ten pairs of zeros give a wavy ramp. A hundred pairs put the corners at the prime powers, with steps of height $\log p$.

\begin{figure}[ht]
\centering
\fig{ch07-explicit}
\caption{The staircase $\psi(x)$ (charcoal) for $1.5\le x\le30$, and the right-hand side of the explicit formula truncated to the first $10$ pairs of zeros (ochre, dashed) and to the first $100$ pairs (coral). With $100$ pairs, on this range, the error stays below $0.34$ at points more than $0.3$ from an integer.}
\label{fig:explicit}
\end{figure}

A zero with real part $\beta$ produces a wave of size about $x^\beta$. The prime number theorem says that $\psi(x)-x$ is small beside $x$. The waves say how small, once one knows how far right the zeros sit. This chapter proves the formula for $\psi_0$, then three consequences. The zero-free region of Chapter~\ref{ch:zeros} gives de la Vallée Poussin's error term. Von Koch's theorem says the Riemann hypothesis is the same statement as $\psi(x)=x+O(\sqrt x\log^2x)$, and the same statement as $\psi(x)=x+O(x^{1/2+\varepsilon})$ for every $\varepsilon>0$. The last section explains the race, promised in Chapter~\ref{ch:dirichlet}, between primes $\equiv3\pmod4$ and primes $\equiv1\pmod4$.

The method is the contour argument of Chapter~\ref{ch:pnt}, with sharper bounds. Write a truncated $\psi_0$ as an integral of $-\zeta'/\zeta$, push the contour left across the strip, and collect a residue at each zero. Chapter~\ref{ch:zeros} supplies the partial-fraction formula for $\zeta'/\zeta$, the count of the zeros, and a zero-free region.

\begin{think}
Look at the formula at a prime $p$. The left side jumps by $\log p$ when $x$ passes $p$. The terms $x$, $\log2\pi$ and $\frac12\log(1-x^{-2})$ are continuous. Each single term $x^\rho/\rho$ is a smooth function of $x$. What does the jump tell you about the way $\sum_\rho x^\rho/\rho$ converges?
\end{think}

\section{Perron's formula}

A sum over integers can be turned into an integral by a switch. For $y>0$ let $\delta(y)=1$ if $y>1$, let $\delta(1)=\frac12$, and let $\delta(y)=0$ if $y<1$. For $c>0$ and $T>0$ set
\[
I(y,T)=\frac1{2\pi i}\int_{c-iT}^{c+iT}\frac{y^s}s\dd s,
\]
where $y^s=e^{s\log y}$ and the path is the vertical segment. The integral is trying to equal $\delta(y)$.

\begin{lemma}[Perron]\label{lem:perron}
Let $y>0$, $c>0$ and $T>0$. If $y\ne1$,
\[
\bigl|I(y,T)-\delta(y)\bigr|\le y^c\min\Bigl(1,\frac1{\pi T|\log y|}\Bigr).
\]
If $y=1$, then $\bigl|I(1,T)-\frac12\bigr|\le\frac c{\pi T}$.
\end{lemma}
\begin{proof}
\emph{Case $y>1$.} Far to the left, $|y^s|=y^{\operatorname{Re}s}$ is small, so the segment moves left. For $U>0$ the function $y^s/s$ is holomorphic except for a simple pole at $0$ with residue $1$. Apply the residue theorem (Volume~5, Theorem~9.26) to the rectangle with vertices $c\pm iT$ and $-U\pm iT$. Then $I(y,T)-1$ equals $-\frac1{2\pi i}$ times the integral over the other three sides, traversed from $c+iT$ to $-U+iT$ to $-U-iT$ to $c-iT$. On a horizontal side, $|s|\ge T$, so the integral of one horizontal side is at most
\[
\frac1T\int_{-\infty}^c y^\sigma\dd\sigma=\frac{y^c}{T\log y}.
\]
The left side has length $2T$ and the integrand is at most $y^{-U}/U$ there, which tends to $0$ as $U\to\infty$. Hence
\[
|I(y,T)-1|\le\frac1{2\pi}\cdot\frac{2y^c}{T\log y}=\frac{y^c}{\pi T\log y}.
\]
For the cruder bound, replace the rectangle by the arc of the circle $|s|=R$, $R=|c+iT|$, from $c+iT$ to $c-iT$ through the left half-plane. The segment and the arc enclose $0$. On the arc, $\operatorname{Re}s\le c$, so $|y^s/s|\le y^c/R$, and the arc has length less than $2\pi R$. So $|I(y,T)-1|\le y^c$.

\emph{Case $y<1$.} Now $|y^s|$ is small far to the right. The same two contours, drawn to the right of the segment, enclose no pole. The same estimates give $|I(y,T)|\le y^c/(\pi T|\log y|)$ and $|I(y,T)|\le y^c$.

\emph{Case $y=1$.} With $s=c+it$,
\[
I(1,T)=\frac1{2\pi}\int_{-T}^T\frac{\mathrm dt}{c+it}=\frac1{2\pi}\int_{-T}^T\frac{c\dd t}{c^2+t^2}=\frac1\pi\arctan\frac Tc.
\]
The odd part $-it/(c^2+t^2)$ integrates to $0$. So
\[
\frac12-I(1,T)=\frac1\pi\Bigl(\frac\pi2-\arctan\frac Tc\Bigr)=\frac1\pi\arctan\frac cT\le\frac c{\pi T}.
\]
\end{proof}

Here is a numerical check of the bound. Take $y=2$, $c=1$ and $T=4$. Writing $s=1+it$ turns the integral into
\[
I(2,4)=\frac1{2\pi}\int_{-4}^{4}\frac{2^{1+it}}{1+it}\dd t=1.192.
\]
The distance to $1$ is $0.192$. Lemma~\ref{lem:perron} allows at most $2/(\pi\cdot4\cdot\log2)=0.230$. The error sits under the bound.

The switch is blurred near $y=1$. Truncating at height $T$ smears the jump over a window of width about $1/T$ in the variable $\log y$. The error terms below are the price of that blur.

Since $\psi$ jumps at prime powers, the value used there is the midpoint. Let
\[
\psi_0(x)=\psi(x)-\frac12\Lambda(x)
\]
when $x$ is an integer, and $\psi_0(x)=\psi(x)$ otherwise. For $x\ge2$, let $\langle x\rangle$ be the distance from $x$ to the nearest prime power other than $x$ itself.

\begin{proposition}[truncated Perron formula]\label{prop:truncperron}
Let $x\ge2$, $T\ge2$ and $c=1+\frac1{\log x}$. Then
\[
\begin{aligned}\psi_0(x)&=\frac1{2\pi i}\int_{c-iT}^{c+iT}\Bigl(-\frac{\zeta'}\zeta(s)\Bigr)\frac{x^s}s\dd s\eqbr+O\Bigl(\frac{x\log^2x}T+\log x\min\Bigl(1,\frac x{T\langle x\rangle}\Bigr)\Bigr).\end{aligned}
\]
The implied constant is absolute.
\end{proposition}
\begin{proof}
On the line $\operatorname{Re}s=c$, the series $-\zeta'/\zeta(s)=\sum_n\Lambda(n)n^{-s}$ converges absolutely and uniformly (Lemma~\ref{lem:logderiv}), so the integral may be taken term by term. It equals $\sum_n\Lambda(n)I(x/n,T)$. Since $\psi_0(x)=\sum_n\Lambda(n)\delta(x/n)$, Lemma~\ref{lem:perron} bounds the difference $R$ by
\[
\begin{aligned}|R|&\le\sum_{n\ne x}\Lambda(n)\Bigl(\frac xn\Bigr)^c\min\Bigl(1,\frac1{\pi T|\log\frac xn|}\Bigr)\eqbr+\frac{c\Lambda(x)}{\pi T},\end{aligned}
\]
with $\Lambda(x)=0$ if $x$ is not an integer. Note $x^c=ex$ and $c\le1+1/\log2<3$. The last term is $O(\log x/T)$. Split the sum.

\emph{Far from $x$: $n\le x/2$ or $n\ge2x$.} Then $|\log(x/n)|\ge\log2$, and these terms contribute at most
\[
\frac{x^c}{\pi T\log2}\sum_n\frac{\Lambda(n)}{n^c}=\frac{ex}{\pi T\log2}\Bigl(-\frac{\zeta'}\zeta(c)\Bigr).
\]
By the second inequality in the proof of Theorem~\ref{thm:zerofree}, $-\zeta'/\zeta(c)\le1/(c-1)+A=\log x+A$ when $c\le2$, and the series $\sum_n\Lambda(n)n^{-2}$ bounds it when $c\ge2$. This part is $O(x\log x/T)$.

\emph{Near $x$: $x/2<n<2x$, $n\ne x$.} Here $(x/n)^c\le2^3=8$ and $\Lambda(n)\le\log(2x)$. Write $n=x(1+u)$ with $-\frac12<u<1$. For $u>0$, $\log(1+u)\ge u/(1+u)\ge u/2$. For $u<0$, $|\log(1+u)|\ge|u|$. So $|\log(x/n)|\ge|n-x|/(2x)$. The integers $n$ with $|n-x|\ge1$ in this range contribute at most
\[
8\log(2x)\sum_{1\le|n-x|<2x}\frac{2x}{\pi T|n-x|}=O\Bigl(\frac{x\log^2x}T\Bigr),
\]
since for each $k\ge1$ at most two integers have $k\le|n-x|<k+1$. At most two integers $n\ne x$ have $|n-x|<1$. If such an $n$ is a prime power, then $|n-x|\ge\langle x\rangle$, and its term is at most $8\log(2x)\min\bigl(1,2x/(\pi T\langle x\rangle)\bigr)$.
\end{proof}

\section{The explicit formula}

The segment now moves left, all the way to $-\infty$. Three pieces have to be controlled: the integrand on the far left, the integrand on the horizontal sides, and the residues. The horizontal sides are the delicate ones. They cross the critical strip, where $\zeta'/\zeta$ has poles. The height is chosen so that the side passes between the zeros, as far from all of them as their density allows. Figure~\ref{fig:perronbox} is the rectangle. The zeros drawn there are a sketch: their heights are not the true ordinates.

\begin{figure}[ht]
\centering
\fig{ch07-new-perron}
\caption{The contour in Perron's formula is the coral segment at $\operatorname{Re}s=c$. The charcoal rectangle is the one used to move that segment left. The point $1$ produces the main term $x$, and the point $0$ produces $-\log2\pi$. The dots in the strip are a sketch of zeros, placed by hand.}
\label{fig:perronbox}
\end{figure}

\begin{lemma}[good heights]\label{lem:goodT}
There is an absolute constant $C$ such that for every $T\ge3$ there is a height $T'\in[T,T+1]$ with
\[
\Bigl|\frac{\zeta'}\zeta(\sigma\pm iT')\Bigr|\le C\log^2T\qquad(-1\le\sigma\le3).
\]
\end{lemma}
\begin{proof}
By Corollary~\ref{cor:nunit}, the number $m$ of ordinates in $[T,T+1]$ is at most $C_0\log T$. They cut $[T,T+1]$ into at most $m+1$ intervals, and one of those has length at least $1/(m+1)$. Let $T'$ be its midpoint. Every ordinate $\gamma$ lies outside the open interval, so $|T'-\gamma|\ge1/(2(m+1))$. For $-1\le\sigma\le2$, Proposition~\ref{prop:nearzeros} writes
\[
\frac{\zeta'}\zeta(\sigma+iT')=\sum_{|\gamma-T'|\le1}\frac1{s-\rho}+O(\log T).
\]
The sum has $O(\log T)$ terms, again by Corollary~\ref{cor:nunit}, and each has modulus at most $1/|T'-\gamma|\le2(m+1)=O(\log T)$. For $2\le\sigma\le3$, $|\zeta'/\zeta(s)|\le\sum_n\Lambda(n)n^{-2}$. The relation $\zeta'/\zeta(\overline s)=\overline{\zeta'/\zeta(s)}$ handles $-T'$.
\end{proof}

On the far left, the functional equation reflects $\zeta'/\zeta$ into the half-plane of absolute convergence. What remains is a logarithm from the Gamma factor.

\begin{lemma}[the left half-plane]\label{lem:leftzeta}
For $\operatorname{Re}s\le-1$, and $|s+2k|\ge1$ for every integer $k\ge1$,
\[
\frac{\zeta'}\zeta(s)=O\bigl(\log(|s|+2)\bigr).
\]
\end{lemma}
\begin{proof}
Take the logarithmic derivative of $\pi^{-s/2}\Gamma(s/2)\zeta(s)=\pi^{-(1-s)/2}\Gamma((1-s)/2)\zeta(1-s)$ (Volume~5, Theorem~10.8):
\[
\begin{aligned}\frac{\zeta'}\zeta(s)&=\log\pi-\frac12\psi_\Gamma\Bigl(\frac s2\Bigr)\eqbr-\frac12\psi_\Gamma\Bigl(\frac{1-s}2\Bigr)-\frac{\zeta'}\zeta(1-s).\end{aligned}
\]
Here $\operatorname{Re}(1-s)\ge2$, so $|\zeta'/\zeta(1-s)|\le\sum_n\Lambda(n)n^{-2}$. The point $(1-s)/2$ has real part at least $1$, so Theorem~\ref{thm:stirlingc}(c), with $\delta=\pi/2$, gives
\[
\psi_\Gamma\Bigl(\frac{1-s}2\Bigr)=\log\frac{1-s}2+O(1)=O\bigl(\log(|s|+2)\bigr).
\]
For $\psi_\Gamma(s/2)$, reflect first. The logarithmic derivative of $\Gamma(z)\Gamma(1-z)=\pi/\sin(\pi z)$ (Volume~5, Theorem~9.36(c)) gives $\psi_\Gamma(z)=\psi_\Gamma(1-z)-\pi\cot(\pi z)$. With $z=s/2$, the point $1-s/2$ has real part at least $3/2$, so $\psi_\Gamma(1-s/2)=O(\log(|s|+2))$ as before. And $z$ stays at distance at least $1/2$ from the integers $\le-1$, and at distance at least $1/2$ from $0$ and the positive integers, since $\operatorname{Re}z\le-1/2$. The function $\cot(\pi z)$ is $1$-periodic, continuous away from the integers, and tends to $\mp i$ as $\operatorname{Im}z\to\pm\infty$, so it is bounded on this set.
\end{proof}

Two constants appear in the residues. First, $\zeta(0)=-\frac12\ne0$ (Volume~5, Corollary~10.9(b)), and
\begin{equation}\label{eq:zetaprime0}
\frac{\zeta'}\zeta(0)=\log2\pi.
\end{equation}
To see it, let $s\to0$ in \eqref{eq:zetapf}. The sum over zeros vanishes at $s=0$, and $-1/(s-1)\to1$. From $\Gamma(z+1)=z\Gamma(z)$ and $\psi_\Gamma(1)=-\gamma$,
\[
\psi_\Gamma(z)=\psi_\Gamma(z+1)-\frac1z=-\frac1z-\gamma+O(z),
\]
so $-1/s-\frac12\psi_\Gamma(s/2)\to\gamma/2$. Hence $\zeta'/\zeta(0)=B+1+\frac12\log\pi+\gamma/2$. With $B=\frac12\log(4\pi)-1-\gamma/2$ (Exercise~\ref{ex:digammahalf}), this is
\[
\frac12\log(4\pi)+\frac12\log\pi=\log(2\pi).
\]
Second, the trivial zeros $-2,-4,\dots$ are simple. They are the poles of $\Gamma(s/2)$, which are simple, and $\xi$ does not vanish at those points.

Before the contour moves, here is the left side at a point you can add by hand. The prime powers up to $6$ are $2,3,4,5$, and $6$ itself is not a prime power. So
\[
\psi_0(6)=\log2+\log3+\log2+\log5=\log60=4.094.
\]
The theorem has to reproduce this number, up to the size of the error at the height $T$ you stop at.

\begin{theorem}[the explicit formula, von Mangoldt 1895]\label{thm:explicit}
For $x\ge2$ and $T\ge3$,
\[
\begin{aligned}\psi_0(x)&=x-\sum_{|\gamma|<T}\frac{x^\rho}\rho-\log2\pi-\frac12\log\Bigl(1-\frac1{x^2}\Bigr)\eqbr+O\Bigl(\frac{x\log^2(xT)}T+\log x\min\Bigl(1,\frac x{T\langle x\rangle}\Bigr)\Bigr),\end{aligned}
\]
the sum running over the non-trivial zeros $\rho=\beta+i\gamma$ with $|\gamma|<T$, counted with multiplicity. The implied constant is absolute.
\end{theorem}
\begin{proof}
\emph{A good height.} Take $T'\in[T,T+1]$ as in Lemma~\ref{lem:goodT}. Proposition~\ref{prop:truncperron} at height $T'$ has the same error term, up to a constant factor. The zeros with $T\le|\gamma|<T'$ number $O(\log T)$ (Corollary~\ref{cor:nunit}), and each satisfies $|x^\rho/\rho|\le x/T$. Replacing the condition $|\gamma|<T$ by $|\gamma|<T'$ costs $O(x\log T/T)$.

\emph{The rectangle.} Let $K\ge1$ be an integer, set $U=2K+1$, and let $\mathcal R$ be the rectangle with vertices $c\pm iT'$ and $-U\pm iT'$, where $c=1+1/\log x$. The function $F(s)=-(\zeta'/\zeta(s))\,x^s/s$ has no pole on $\partial\mathcal R$: there are no zeros on $\operatorname{Re}s=c>1$, the height $T'$ is not an ordinate, and $-U$ is odd. Inside, the poles and residues are:
\begin{itemize}
\item at $s=1$, where $-\zeta'/\zeta$ has a simple pole of residue $1$: residue $x$;
\item at each non-trivial zero $\rho$ of multiplicity $m$, where $-\zeta'/\zeta$ has residue $-m$: residue $-m\,x^\rho/\rho$;
\item at $s=0$, where $-\zeta'/\zeta$ is holomorphic: residue $-\zeta'/\zeta(0)=-\log2\pi$;
\item at $s=-2k$ for $1\le k\le K$: residue $-x^{-2k}/(-2k)=x^{-2k}/(2k)$.
\end{itemize}
By the residue theorem, $\frac1{2\pi i}\int_{c-iT'}^{c+iT'}F$ equals the sum of these residues minus $\frac1{2\pi i}$ times the integral of $F$ over the other three sides.

\emph{The other sides.} On the left side, $|F(s)|\le C\log(U+T')\,x^{-U}/U$ by Lemma~\ref{lem:leftzeta}, and the side has length $2T'$. Its integral tends to $0$ as $K\to\infty$. On the horizontal sides, $|s|\ge T'$. For $-1\le\sigma\le c$, Lemma~\ref{lem:goodT} gives $|F(s)|\le C\log^2T\cdot x^\sigma/T$, and
\[
\int_{-1}^c x^\sigma\dd\sigma\le\frac{x^c}{\log x}=\frac{ex}{\log x}.
\]
For $\sigma\le-1$, Lemma~\ref{lem:leftzeta} applies (the side sits at height $T'\ge3$), so $|F(s)|\le C\log(|\sigma|+T+3)\,x^\sigma/T$, and the integral over $\sigma\le-1$ is $O(\log T/(xT))$. Altogether the horizontal sides contribute $O(x\log^2T/(T\log x))$.

\emph{Conclusion.} Let $K\to\infty$. The trivial zeros contribute
\[
\sum_{k\ge1}\frac{x^{-2k}}{2k}=-\frac12\log(1-x^{-2}).
\]
Collecting the errors,
\[
\frac{x\log^2x}T+\frac{x\log T}T+\frac{x\log^2T}{T\log x}=O\Bigl(\frac{x\log^2(xT)}T\Bigr).
\]
\end{proof}

Let $T\to\infty$ with $x$ fixed. The error tends to $0$, because $\langle x\rangle>0$. This is Corollary~\ref{cor:exactexplicit}, the formula from the introduction.

\begin{corollary}[the exact formula]\label{cor:exactexplicit}
For $x\ge2$,
\[
\begin{aligned}\psi_0(x)&=x-\lim_{T\to\infty}\sum_{|\gamma|<T}\frac{x^\rho}\rho\eqbr-\log2\pi-\frac12\log\Bigl(1-\frac1{x^2}\Bigr).\end{aligned}
\]
\end{corollary}

At $x=10$ the right-hand side, stopped after $100$ pairs of zeros, equals $7.826$. The left side is $\psi_0(10)=7.832$. Ten pairs give $7.787$. The symmetric limit is doing the work; a finite list of zeros is an approximation.

The series $\sum x^\rho/\rho$ does not converge absolutely: $\sum 1/|\rho|$ diverges (Exercise~\ref{ex:zerosums}). It converges in the symmetric sense of the limit, and the convergence is not uniform near a prime power. There a sum of smooth functions builds a jump, in the same way that a Fourier series builds a jump. The overshoot is visible in Figure~\ref{fig:explicit}. Exercise~\ref{ex:cesaro} replaces the sharp cut by a sloping weight, and for that weight the series converges absolutely.

Figure~\ref{fig:threewaves} is a picture of that idea, drawn with three sines and a step. It is the beginning of the Fourier series for a square wave. It is not a partial sum over zeros of $\zeta$.

\begin{figure}[ht]
\centering
\fig{ch07-new-waves}
\caption{Three sine waves, $\frac4\pi\sin u$, $\frac4{3\pi}\sin3u$ and $\frac4{5\pi}\sin5u$, and their sum (coral), beside the square wave they begin to follow (charcoal, dashed). At $u=\pi/2$ the sum equals $\frac4\pi\bigl(1-\frac13+\frac15\bigr)=\frac{52}{15\pi}\approx1.103$, a little above the step height $1$. This is a picture of smooth waves adding toward a jump. The frequencies are $1,3,5$, not ordinates of zeros of $\zeta$.}
\label{fig:threewaves}
\end{figure}

\begin{remark}[von Mangoldt, 1895]
Hans von Mangoldt proved the explicit formula in 1895, in the \emph{Journal f\"ur die reine und angewandte Mathematik}, volume 114, pages 255--305. Riemann's memoir of 1859 had stated a formula for a different staircase and had not given a complete proof. The identity von Mangoldt proved, in the form recorded here, is Corollary~\ref{cor:exactexplicit}: the left side is $\psi_0(x)$, the value halfway up a jump, and the sum over zeros is a symmetric limit. The paper is listed in the sources at the end of the chapter.

\portrait{mangoldt}{Hans von Mangoldt (1854--1925). Photograph from the Weierstrass album, between 1870 and 1885. Public domain.}
\end{remark}

\begin{remark}[Riemann's formula, quoted]
Riemann's own formula was for $\pi(x)$, and it is more intricate. Let $J(x)=\sum_{p^k\le x}1/k$, which counts prime powers with weight $1/k$, with the midpoint convention at the jumps. Then $\pi(x)=\sum_{n\ge1}\mu(n)/n\,J(x^{1/n})$ by M\"obius inversion, and Riemann stated
\[
\begin{aligned}J(x)&=\operatorname{li}(x)-\sum_\rho\operatorname{li}(x^\rho)-\log2\eqbr+\int_x^\infty\frac{\mathrm dt}{t(t^2-1)\log t},\end{aligned}
\]
where $\operatorname{li}$ is the logarithmic integral, continued to complex arguments. Von Mangoldt proved this version as well. It is not used below. Edwards's book, Chapters~1 and~3, tells the story and gives the proof. The $\psi$ version is simpler because $\zeta'/\zeta$ has simple poles with integer residues, whereas $J$ corresponds to $\log\zeta$, which has logarithmic singularities.
\end{remark}

\section{The error term in the prime number theorem}

The formula says that $\psi(x)-x$ is a sum of waves $x^\rho/\rho$, each of size $x^\beta/|\rho|$. A zero-free region bounds every $\beta$, and so bounds every wave. There are infinitely many waves, and the region of Theorem~\ref{thm:zerofree} gets narrower as $\gamma$ grows. Truncating at a height $T$ balances the tail against the waves that were kept.

\begin{lemma}\label{lem:sumrho}
For $T\ge3$, $\displaystyle\sum_{|\gamma|<T}\frac1{|\rho|}=O(\log^2T)$.
\end{lemma}
\begin{proof}
No zero is $0$, since $\zeta(0)\ne0$, and only finitely many have $|\gamma|<1$: they contribute $O(1)$. By Corollary~\ref{cor:nunit} and the symmetry of the zeros, the zeros with $k\le|\gamma|<k+1$ number $O(\log(k+2))$, and each has $|\rho|\ge k$. The sum is
\[
O\Bigl(1+\sum_{k<T}\frac{\log(k+2)}k\Bigr)=O(\log^2T).
\]
\end{proof}

\begin{theorem}[de la Vall\'ee Poussin, 1899]\label{thm:pnterror}
There is a constant $c>0$ such that
\[
\psi(x)=x+O\bigl(xe^{-c\sqrt{\log x}}\bigr).
\]
\end{theorem}
\begin{proof}
Let $x$ be large. Since $|\psi(x)-\psi_0(x)|\le\log x$, and the minimum in Theorem~\ref{thm:explicit} is at most $1$,
\[
|\psi(x)-x|\le\sum_{|\gamma|<T}\frac{x^\beta}{|\rho|}+O\Bigl(\frac{x\log^2(xT)}T+\log x\Bigr).
\]
By Theorem~\ref{thm:zerofree}, every zero with $|\gamma|<T$ has
\[
\beta\le1-\frac{c_1}{\log(|\gamma|+2)}\le1-\frac{c_1}{\log(T+2)},
\]
so $x^\beta\le x\exp\bigl(-c_1\log x/\log(T+2)\bigr)$. Lemma~\ref{lem:sumrho} then gives
\[
\begin{aligned}|\psi(x)-x|&=O\Bigl(x\log^2T\,e^{-c_1\frac{\log x}{\log(T+2)}}\eqbr+\frac{x\log^2(xT)}T+\log x\Bigr).\end{aligned}
\]
The exponent uses $\log(T+2)$, the same denominator as the bound on $\beta$. Choose $T$ to balance the two main terms: $\log T=\sqrt{\log x}$. Then $\log T/\log(T+2)\to1$. For large $x$ the ratio is at least $\frac12$, so
\[
\frac{\log x}{\log(T+2)}=\sqrt{\log x}\cdot\frac{\log T}{\log(T+2)}\ge\frac12\sqrt{\log x},
\]
and the first term is $O\bigl(x\log x\,e^{-(c_1/2)\sqrt{\log x}}\bigr)$. The second term is $O\bigl(x\log^2x\,e^{-\sqrt{\log x}}\bigr)$. Any $c<\min(c_1/2,1)$ works.
\end{proof}

The error $xe^{-c\sqrt{\log x}}$ is smaller than $x/\log^kx$ for every $k$, and larger than $x^{1-\varepsilon}$ for every $\varepsilon>0$. A power saving would need a zero-free strip $\beta\le1-\eta$ of fixed width. No such strip is known.

\begin{corollary}\label{cor:pnterrorpi}
With a smaller $c>0$, $\pi(x)=\operatorname{Li}(x)+O\bigl(xe^{-c\sqrt{\log x}}\bigr)$.
\end{corollary}
\begin{proof}
By Lemma~\ref{lem:abelint} with $a_p=\log p$, $f(t)=1/\log t$ and $y=2$ (and $\theta(2)=\log2$),
\[
\pi(x)=\frac{\theta(x)}{\log x}+\int_2^x\frac{\theta(t)}{t\log^2t}\dd t.
\]
Integrating by parts,
\[
\operatorname{Li}(x)=\frac x{\log x}-\frac2{\log2}+\int_2^x\frac{\mathrm dt}{\log^2t}.
\]
Subtracting,
\[
\begin{aligned}\pi(x)-\operatorname{Li}(x)&=\frac{\theta(x)-x}{\log x}+\frac2{\log2}\eqbr+\int_2^x\frac{\theta(t)-t}{t\log^2t}\dd t.\end{aligned}
\]
Since $\psi(x)-\theta(x)=O(\sqrt x\log x)$ (Chapter~\ref{ch:mertens}), one has $\theta(t)-t=O\bigl(te^{-c\sqrt{\log t}}\bigr)$. The integral is then at most a constant times
\[
\int_2^{\sqrt x}\mathrm dt+\int_{\sqrt x}^xe^{-c\sqrt{\log t}}\dd t\le\sqrt x+xe^{-c\sqrt{(\log x)/2}}.
\]
\end{proof}

\section{The Riemann hypothesis as a statement about primes}

If every non-trivial zero has real part $\frac12$, every wave has amplitude $\sqrt x/|\rho|$, and the waves together stay within $\sqrt x$ times a power of $\log x$. A single zero with $\beta>\frac12$ would produce a wave of size $x^\beta$, and the other waves could not cancel it for all large $x$. So a statement about where $\zeta$ vanishes is the same statement as a bound on $\psi(x)-x$. The shape proved here is $O(\sqrt x\log^2x)$, and also $O(x^{1/2+\varepsilon})$ for every $\varepsilon>0$. The argument does not reach a smaller power of $\log x$.

\begin{theorem}[von Koch, 1901]\label{thm:rhequiv}
The following are equivalent:
\begin{enumerate}[label=(\alph*)]
\item the Riemann hypothesis: every non-trivial zero has real part $\frac12$;
\item $\psi(x)=x+O(\sqrt x\log^2x)$;
\item for every $\varepsilon>0$, $\psi(x)=x+O(x^{1/2+\varepsilon})$.
\end{enumerate}
\end{theorem}
\begin{proof}
(a)$\Rightarrow$(b). Take $T=x$ in Theorem~\ref{thm:explicit}. Each $|x^\rho|=\sqrt x$, so by Lemma~\ref{lem:sumrho} the sum over zeros is $O(\sqrt x\log^2x)$. The error term is
\[
O\Bigl(\frac{x\log^2(x^2)}x+\log x\Bigr)=O(\log^2x).
\]
Also $|\psi-\psi_0|\le\log x$.

(b)$\Rightarrow$(c) is immediate: $\sqrt x\log^2x=O(x^{1/2+\varepsilon})$ for every $\varepsilon>0$.

(c)$\Rightarrow$(a). For $\operatorname{Re}s>1$, Lemma~\ref{lem:abelint} with $a_n=\Lambda(n)$ and $f(t)=t^{-s}$, then $X\to\infty$ (using $\psi(X)X^{-s}\to0$ since $\psi(X)=O(X)$), gives
\[
\begin{aligned}-\frac{\zeta'}\zeta(s)&=s\int_1^\infty\frac{\psi(t)}{t^{s+1}}\dd t\\&=\frac s{s-1}+s\int_1^\infty\frac{\psi(t)-t}{t^{s+1}}\dd t.\end{aligned}
\]
Fix $\varepsilon>0$. Under (c), $|\psi(t)-t|\le Ct^{1/2+\varepsilon}$, so the last integral converges absolutely and uniformly on the half-plane $\operatorname{Re}s\ge\frac12+2\varepsilon$, and on every half-plane $\operatorname{Re}s\ge\frac12+2\varepsilon+\eta$ with $\eta>0$. So the integral is holomorphic on $\operatorname{Re}s>\frac12+2\varepsilon$ (Volume~5, Theorem~9.15), and the right-hand side $G(s)$ is holomorphic on $H=\{\operatorname{Re}s>\frac12+2\varepsilon\}$ except at $1$. Let $H'$ be $H$ with $1$ and the zeros of $\zeta$ removed. The set $H'$ is open and connected, because a discrete set was removed. On $H'$, both $-\zeta'/\zeta$ and $G$ are holomorphic, and they agree for $\operatorname{Re}s>1$. By the identity theorem they agree on $H'$. If $\zeta(\rho)=0$ with $\rho\in H$, then $-\zeta'/\zeta$ has a pole at $\rho$, while $G$ is bounded near $\rho$. So no such zero exists. Every zero has $\beta\le\frac12+2\varepsilon$. This holds for every $\varepsilon>0$, so $\beta\le\frac12$. The symmetry $\rho\mapsto1-\overline\rho$ (Proposition~\ref{prop:xi}) gives $\beta\ge\frac12$ as well.
\end{proof}

\begin{think}
The proof of (c)$\Rightarrow$(a) uses (c) only to make an integral converge. Run it with the hypothesis $\psi(x)-x=O(x^\theta)$ for some $\theta<\frac12$. What do you get, and why is that absurd? (Exercise~\ref{ex:omegahalf}.) So if the Riemann hypothesis holds, $\sqrt x$ is the true size of the error up to powers of $\log x$. And whether the hypothesis holds or not, no error term can be $O(x^\theta)$ for a $\theta<\frac12$.
\end{think}

\begin{remark}[an explicit constant, quoted]
Schoenfeld (1976) made (a)$\Rightarrow$(b) explicit: if the Riemann hypothesis holds, then $|\psi(x)-x|<\frac1{8\pi}\sqrt x\log^2x$ for all $x\ge73.2$. The constant $\frac1{8\pi}$ equals $0.0398$. On the range $73.2\le x\le10^7$ the ratio $|\psi(x)-x|/(\sqrt x\log^2x)$ reaches $0.0365$ and no higher, so the inequality holds on that range by direct computation. On $10^3\le x\le10^7$ the ratio $|\psi(x)-x|/\sqrt x$ reaches $0.803$ and no higher. Littlewood proved in 1914 that $\psi(x)-x$ is infinitely often larger than a constant times $\sqrt x\log\log\log x$, and infinitely often smaller than the negative of that quantity. Theorem~\ref{thm:rhequiv} stops at $O(\sqrt x\log^2x)$, which is the shape its proof gives.
\end{remark}

\section{Chebyshev's bias}

In 1853 Chebyshev wrote to Fuss that there seem to be more primes of the form $4n+3$ than of the form $4n+1$. Figure~\ref{fig:race} of Chapter~\ref{ch:dirichlet} shows the count. Up to $26860$, the class $3$ is ahead or tied at every integer $x$. At $x=26861$, which is itself a prime $\equiv1\pmod4$, the class $1$ goes ahead by one prime. The next prime is $26863\equiv3\pmod4$, and the race is tied again. Leech found the crossing in 1957. Up to $10^7$, the class $1$ is strictly ahead at $3406$ integers. The prime number theorem for progressions says the two classes are the same size in the limit. The identity below says where the difference of size $\sqrt x$ comes from.

Every odd square is $\equiv1\pmod4$. The function $\psi(x;4,1)$ counts prime powers, so it receives the squares of the odd primes. The counts of primes alone do not. The weighted difference with no main term is
\[
\psi(x,\chi_{-4})=\sum_{n\le x}\chi_{-4}(n)\Lambda(n),
\]
and the squares sit inside that sum. A race of prime powers is the fair race. A race of primes alone has given the squares to the class $1$.

Let $\theta(x;q,a)=\sum_{p\le x,\,p\equiv a\,(q)}\log p$.

\begin{proposition}\label{prop:biasidentity}
For $x\ge2$, with $\chi=\chi_{-4}$,
\[
\begin{aligned}\theta(x;4,3)-\theta(x;4,1)&=\theta(\sqrt x)-\psi(x,\chi)\eqbr+O(x^{1/3}\log x).\end{aligned}
\]
In particular
\[
\frac{\theta(x;4,3)-\theta(x;4,1)}{\sqrt x}=1-\frac{\psi(x,\chi)}{\sqrt x}+o(1).
\]
\end{proposition}
\begin{proof}
Write $\psi(x,\chi)=\sum_{p^k\le x}\chi(p^k)\log p$ and group the terms by the exponent $k$. For $k=1$ the sum is $\sum_{p\le x}\chi(p)\log p=\theta(x;4,1)-\theta(x;4,3)$. For $k=2$, one has $\chi(p^2)=\chi(p)^2=1$ for odd $p$ and $\chi(4)=0$, so the squares contribute $\theta(\sqrt x)-\log2$ when $x\ge4$. For $k\ge3$ the terms add up to at most
\[
\sum_{3\le k\le\log_2x}\theta(x^{1/k})\le\log_2x\cdot\theta(x^{1/3})=O(x^{1/3}\log x),
\]
by Proposition~\ref{prop:chebyshev}. Rearranging gives the identity. Finally $\theta(\sqrt x)\sim\sqrt x$ by the prime number theorem.
\end{proof}

The identity is unconditional. It compares two weighted counts of primes. The term $\theta(\sqrt x)$ is a head start of size $\sqrt x$ for the class $3$. Who is ahead at a given $x$ depends on $\psi(x,\chi)$ as well, and the identity does not estimate $\psi(x,\chi)$. The passage from these weighted counts $\theta$ to the plain counts $\pi(x;4,3)$ and $\pi(x;4,1)$ is a further step. The next remark reads the identity as a sum of waves. That reading uses a functional equation this volume does not prove, and no theorem of the chapter depends on it.

\begin{remark}[the waves of $L(s,\chi_{-4})$; the functional equation of $L(s,\chi)$ is admitted here, proved in Volume~13, Chapter~12]
The method of this chapter applies to $L(s,\chi)$ for a primitive character, once the functional equation of $L(s,\chi)$ is available in place of the functional equation of $\zeta$. That functional equation is \emph{admitted} in this volume. It is proved in Volume~13, Chapter~12, as a corollary of Tate's thesis for $\Q$. Davenport's book proves it in Chapter~9 and draws the explicit formula from it in Chapter~19. Proposition~\ref{prop:biasidentity} does not use it. Since $L(s,\chi_{-4})$ has no pole, the admitted formula has no main term $x$. For $x\ge2$,
\[
\psi_0(x,\chi_{-4})=-\lim_{T\to\infty}\sum_{|\gamma|<T}\frac{x^\rho}\rho+O(1),
\]
the sum running over the non-trivial zeros of $L(s,\chi_{-4})$. The first three positive ordinates are $6.0209$, $10.2438$ and $12.9881$. If every one of those zeros has real part $\frac12$, then with $x=e^u$ the identity of Proposition~\ref{prop:biasidentity} becomes
\[
\frac{\theta(e^u;4,3)-\theta(e^u;4,1)}{e^{u/2}}=1+\sum_{\gamma\neq0}\frac{e^{i\gamma u}}{\frac12+i\gamma}+o(1),
\]
and the sum equals $\sum_{\gamma>0}2\operatorname{Re}(e^{i\gamma u}/\rho)$. The constant $1$ is the head start $\theta(\sqrt x)/\sqrt x$. The sum over zeros is not estimated in this chapter.
\end{remark}

Rubinstein and Sarnak (1994) assumed the generalized Riemann hypothesis for the characters in the race, and they assumed that the positive ordinates are linearly independent over $\Q$. Under those two hypotheses they proved that the set of $x$ where $\pi(x;4,3)>\pi(x;4,1)$ has logarithmic density $0.9959\ldots$: in the variable $u=\log x$, the class $3$ leads $99.59\%$ of the time. That density is their theorem. It is not a consequence of the unconditional identity in Proposition~\ref{prop:biasidentity}, which controls $\theta$ and does not by itself produce a density for $\pi$. Littlewood proved in 1914, with no hypothesis on the zeros, that the lead in the race of $\pi(x;4,3)$ against $\pi(x;4,1)$ changes infinitely often. The same bias, a head start for the non-squares, appears for every modulus $q$ (Exercise~\ref{ex:bias}).

\begin{think}
Primes $\equiv1\pmod3$ and $\equiv2\pmod3$ also race. Squares are $\equiv1\pmod3$. Which class receives the head start of size $\sqrt x$ in the identity of Exercise~\ref{ex:bias}? Check the first few thousand primes, then compare with the identity.
\end{think}

\section{What we proved, and what comes next}

Corollary~\ref{cor:exactexplicit} writes $\psi_0(x)$ as $x$ minus a symmetric sum of $x^\rho/\rho$, minus $\log2\pi$, minus $\frac12\log(1-x^{-2})$. Theorem~\ref{thm:explicit} keeps an explicit error when the sum stops at height $T$. From the zero-free region, Theorem~\ref{thm:pnterror} gives $\psi(x)=x+O(xe^{-c\sqrt{\log x}})$. Theorem~\ref{thm:rhequiv} says the Riemann hypothesis is equivalent to $\psi(x)=x+O(\sqrt x\log^2x)$ and to $\psi(x)=x+O(x^{1/2+\varepsilon})$ for every $\varepsilon>0$. Proposition~\ref{prop:biasidentity} is the unconditional identity behind Chebyshev's race: the class $3$ has a head start $\theta(\sqrt x)$, because the odd squares lie in the class $1$.

The waves bound the size of $\psi(x)-x$. They do not produce primes in intervals shorter than $\sqrt x$, they do not bound the twin primes, and they say nothing about the number of prime factors of a typical integer. Each of those questions asks for integers that avoid one or two residue classes modulo every small prime. A sieve does that counting with no information about the zeros. Chapter~\ref{ch:sieves} builds one.

\section{Exercises}

\begin{exercise}[$\star$]\label{ex:sumgamma}
Show that $\displaystyle\sum_{0<\gamma\le T}\frac1\gamma=\frac1{4\pi}\log^2T+O(\log T)$.
\end{exercise}

\begin{exercise}[$\star$]\label{ex:omegahalf}
Show that $\psi(x)-x$ is not $O(x^\theta)$ for any $\theta<\frac12$.
\end{exercise}

\begin{exercise}[$\star$]\label{ex:mertensRH}
Let $M(x)=\sum_{n\le x}\mu(n)$. Show that if $M(x)=O(x^{1/2+\varepsilon})$ for every $\varepsilon>0$, then the Riemann hypothesis holds. The converse is quoted in the remark after Theorem~\ref{thm:mobiussum}; it is not asked here.
\end{exercise}

\begin{exercise}[$\star\star$]\label{ex:shortprime}
(a) Show that for every $A>0$, the interval $\bigl]x,x+x/\log^Ax\bigr]$ contains $\sim x/\log^{A+1}x$ primes as $x\to\infty$.
(b) Show that the Riemann hypothesis implies that $]x,x+\sqrt x\log^3x]$ contains a prime for all large $x$. (Cram\'er, 1920.)
\end{exercise}

\begin{exercise}[$\star\star$]\label{ex:vonkoch}
Show that the Riemann hypothesis implies $\pi(x)=\operatorname{Li}(x)+O(\sqrt x\log x)$.
\end{exercise}

\begin{exercise}[$\star\star$, the signature exercise]\label{ex:bias}
(a) Show that $\theta(x;3,2)-\theta(x;3,1)=\theta(\sqrt x)-\psi(x,\chi_{-3})+O(x^{1/3}\log x)$.
(b) Modulo $8$ there are four classes of odd numbers, $1,3,5,7$, and every odd square is $\equiv1\pmod8$. Using the four characters modulo $8$, show that for $a\in\{1,3,5,7\}$,
\[
\begin{aligned}\theta(x;8,a)&=\frac{\theta(x)}4-\frac{\varepsilon_a}4\theta(\sqrt x)\\&\quad+\frac14\sum_{\chi\ne\chi_0}\chi(a)\psi(x,\chi)\eqbr+O(x^{1/3}\log x),\end{aligned}
\]
with $\varepsilon_1=3$ and $\varepsilon_a=-1$ for $a\ne1$. Which class carries the coefficient $-\frac34$ in front of $\theta(\sqrt x)$?
\end{exercise}

\begin{exercise}[$\star\star\star$]\label{ex:cesaro}
Show that for $x\ge2$,
\[
\begin{aligned}\sum_{n\le x}\Lambda(n)\Bigl(1-\frac nx\Bigr)&=\frac x2-\sum_\rho\frac{x^\rho}{\rho(\rho+1)}\eqbr-\log2\pi+\frac{\zeta'(-1)}{\zeta(-1)\,x}\\&\quad-\sum_{k\ge1}\frac{x^{-2k}}{2k(2k-1)},\end{aligned}
\]
where now the sum over zeros converges absolutely. (At $x=10$ the left side equals $3.37642$. The right side, with $200$ pairs of zeros, equals $3.37645$. Here $\zeta'(-1)/\zeta(-1)=1.98505$.)
\end{exercise}

\begin{exercise}[$\star$]\label{ex:halfjump}
At the prime $p=5$, compute $\psi(5)$ and $\psi_0(5)$ from the definition, and check that the difference is $\frac12\log5$. Then evaluate
\[
5-\log(2\pi)-\frac12\log\Bigl(1-\frac1{25}\Bigr)
\]
and deduce the value of the symmetric sum $\lim_{T\to\infty}\sum_{|\gamma|<T}5^\rho/\rho$.
\end{exercise}

\begin{exercise}[$\star$]\label{ex:trivten}
Show that $\displaystyle\sum_{k\ge1}\frac{x^{-2k}}{2k}=-\frac12\log\bigl(1-x^{-2}\bigr)$ for $x>1$. At $x=10$, compute the sum to three significant figures. Set the size beside the main term $x$ in Corollary~\ref{cor:exactexplicit}.
\end{exercise}

\begin{exercise}[$\star\star$]\label{ex:balanceT}
In the proof of Theorem~\ref{thm:pnterror} one chooses $\log T=\sqrt{\log x}$.
\begin{enumerate}[label=(\alph*)]
\item Show that $\log T/\log(T+2)\to1$ as $x\to\infty$.
\item Take $\log x=100$, so $T=e^{10}$. Compute $\log T/\log(T+2)$ to three decimal places, and check that it is greater than $\frac12$.
\item With that choice, $\log x/\log(T+2)$ equals $\sqrt{\log x}$ times the ratio in (b). Write the numerical value of $\log x/\log(T+2)$ for $\log x=100$.
\end{enumerate}
\end{exercise}

\begin{exercise}[$\star\star$]\label{ex:psi0koch}
Show that $|\psi(x)-\psi_0(x)|\le\frac12\log x$ for $x\ge2$. Deduce that the three statements in Theorem~\ref{thm:rhequiv} stay equivalent if $\psi$ is replaced by $\psi_0$ throughout. The shape $O(\sqrt x\log^2x)$ does not become a smaller power of $\log x$ under that replacement.
\end{exercise}

\section{Hints}

\begin{hint}{ex:sumgamma}
Write $\frac1\gamma=\frac1T+\int_\gamma^T\frac{\mathrm dt}{t^2}$ and sum: $\sum\frac1\gamma=\frac{N(T)}T+\int\frac{N(t)}{t^2}\dd t$. Then use Theorem~\ref{thm:rvm}.
\end{hint}
\begin{hint}{ex:omegahalf}
Follow (c)$\Rightarrow$(a) in Theorem~\ref{thm:rhequiv}. Volume~5 found a zero at $\frac12+14.1347\ldots i$.
\end{hint}
\begin{hint}{ex:mertensRH}
For $\operatorname{Re}s>1$, $\frac1{\zeta(s)}$ is $\sum\mu(n)n^{-s}$, which is $s\int_1^\infty M(t)t^{-s-1}\dd t$.
\end{hint}
\begin{hint}{ex:shortprime}
Apply Theorem~\ref{thm:pnterror} (respectively Theorem~\ref{thm:rhequiv}) at $x+h$ and at $x$, and subtract. Remove the prime powers $p^k$, $k\ge2$, with Proposition~\ref{prop:chebyshev}.
\end{hint}
\begin{hint}{ex:vonkoch}
Use the identity for $\pi(x)-\operatorname{Li}(x)$ in the proof of Corollary~\ref{cor:pnterrorpi}.
\end{hint}
\begin{hint}{ex:bias}
(a) Copy the proof of Proposition~\ref{prop:biasidentity}; the prime $3$ plays the role of $2$. (b) Let $\theta(x,\chi)=\sum_{p\le x}\chi(p)\log p$. Orthogonality gives $\theta(x;8,a)=\frac14\sum_\chi\chi(a)\theta(x,\chi)$, and $\sum_{\chi\ne\chi_0}\chi(a)=4\cdot\mathbf1_{a\equiv1}-1$. The characters modulo $8$ are real.
\end{hint}
\begin{hint}{ex:cesaro}
First show that for $c>0$ the integral $\frac1{2\pi i}\int_{c-i\infty}^{c+i\infty}\frac{y^s}{s(s+1)}\dd s$ equals $(1-1/y)\mathbf1_{y\ge1}$. Then shift the contour as in Theorem~\ref{thm:explicit}. The factor $1/(s(s+1))$ makes the integral converge.
\end{hint}
\begin{hint}{ex:halfjump}
The prime powers up to $5$ are $2,3,4,5$. Then $\psi_0(5)=\psi(5)-\frac12\log5$, and Corollary~\ref{cor:exactexplicit} solves for the sum over zeros.
\end{hint}
\begin{hint}{ex:trivten}
The series $\sum_{k\ge1}z^k/k=-\log(1-z)$ at $z=x^{-2}$, divided by $2$.
\end{hint}
\begin{hint}{ex:balanceT}
Write $\log(T+2)=\log T+\log(1+2/T)$ and use $\log(1+2/T)\sim2/T$.
\end{hint}
\begin{hint}{ex:psi0koch}
The difference $\psi-\psi_0$ equals $\frac12\Lambda(x)$ at an integer and equals $0$ otherwise. Compare $\log x$ with $\sqrt x\log^2x$ and with $x^{1/2+\varepsilon}$.
\end{hint}

\section{Solutions}

\begin{solution}{ex:sumgamma}
\idea A sum over zeros weighted by a smooth function is an integral against $\mathrm dN(t)$. The main term of $N$ does the work, and its error term costs $O(\log T)$.
\writeup Finitely many ordinates lie in $]0,3[$; they contribute $O(1)$. For $3\le\gamma\le T$, write $\frac1\gamma=\frac1T+\int_\gamma^T\mathrm dt/t^2$. Summing the finite sum,
\[
\begin{aligned}\sum_{3\le\gamma\le T}\frac1\gamma&=\frac{N(T)-N(3^-)}T\eqbr+\int_3^T\frac{N(t)-N(3^-)}{t^2}\dd t.\end{aligned}
\]
By Theorem~\ref{thm:rvm}, $N(T)/T=O(\log T)$ and $N(t)=\frac t{2\pi}\log\frac t{2\pi e}+O(\log t)$. The error contributes $\int_3^\infty(\log t)/t^2\dd t=O(1)$, and the main term is
\[
\begin{aligned}\int_3^T\frac{\log(t/(2\pi e))}{2\pi t}\dd t&=\frac1{4\pi}\Bigl[\log^2\frac t{2\pi e}\Bigr]_3^T\\&=\frac1{4\pi}\log^2T+O(\log T).\end{aligned}
\]
\end{solution}

\begin{solution}{ex:omegahalf}
\idea A power saving in the error term pushes the analytic continuation of $\zeta'/\zeta$ to the left, and the continuation runs into a zero.
\writeup Suppose $\psi(x)-x=O(x^\theta)$ with $\theta<\frac12$. As in the proof of Theorem~\ref{thm:rhequiv}, the function
\[
\frac s{s-1}+s\int_1^\infty(\psi(t)-t)\,t^{-s-1}\dd t
\]
is holomorphic on $\operatorname{Re}s>\theta$ except at $1$, and it equals $-\zeta'/\zeta$ wherever both are holomorphic. But $\zeta$ vanishes at $\rho_1=\frac12+14.1347\ldots i$ (Volume~5 located it by a sign change of Hardy's $Z$), and $\operatorname{Re}\rho_1=\frac12>\theta$, so $-\zeta'/\zeta$ has a pole there while the integral stays bounded nearby. That is a contradiction. So for every $\theta<\frac12$, one has $|\psi(x)-x|>x^\theta$ for arbitrarily large $x$.
\end{solution}

\begin{solution}{ex:mertensRH}
\idea The same mechanism as Theorem~\ref{thm:rhequiv}, with $1/\zeta$ in place of $-\zeta'/\zeta$: a zero of $\zeta$ is a pole of $1/\zeta$.
\writeup For $\operatorname{Re}s>1$, Lemma~\ref{lem:abelint} with $a_n=\mu(n)$, $f(t)=t^{-s}$, and $M(X)X^{-s}\to0$, gives $1/\zeta(s)=s\int_1^\infty M(t)t^{-s-1}\dd t$. Euler's product gives $1/\zeta=\sum\mu(n)n^{-s}$ (Volume~5, Corollary~10.4). If $M(t)=O(t^{1/2+\varepsilon})$ for every $\varepsilon>0$, then for every $\varepsilon>0$ the integral is holomorphic on $\operatorname{Re}s>\frac12+\varepsilon$ (Volume~5, Theorem~9.15), so $1/\zeta$ extends holomorphically there, by the identity theorem on the connected set where $\zeta\ne0$. A zero of $\zeta$ with $\beta>\frac12+\varepsilon$ would be a pole of this extension. Hence $\beta\le\frac12+\varepsilon$ for every $\varepsilon>0$, so $\beta\le\frac12$, and $\beta\ge\frac12$ by symmetry. One fixed $\varepsilon$ gives only $\beta\le\frac12+\varepsilon$.
\end{solution}

\begin{solution}{ex:shortprime}
\idea The number of primes in an interval is the difference of two values of $\psi$, divided by about $\log x$. The count is accurate when the length $h$ beats the error term.
\writeup (a) Let $h=x/\log^Ax$. By Theorem~\ref{thm:pnterror},
\[
\psi(x+h)-\psi(x)=h+O\bigl(xe^{-c\sqrt{\log x}}\bigr)=h(1+o(1)),
\]
since $e^{-c\sqrt{\log x}}\log^Ax\to0$. The prime powers $p^k\le2x$ with $k\ge2$ contribute at most $\psi(2x)-\theta(2x)=O(\sqrt x\log x)=o(h)$, so $\theta(x+h)-\theta(x)\sim h$. Each prime in $]x,x+h]$ has $\log x<\log p\le\log(2x)$, so their number lies between $(\theta(x+h)-\theta(x))/\log(2x)$ and $(\theta(x+h)-\theta(x))/\log x$, and both are $\sim h/\log x$.

(b) Let $h=\sqrt x\log^3x$. Under the Riemann hypothesis, Theorem~\ref{thm:rhequiv} gives $\psi(x+h)-\psi(x)=h+O(\sqrt x\log^2x)$, which is positive for large $x$. For the prime powers with $k\ge2$ in $]x,x+h]$: there are at most $h/(2\sqrt x)+1$ squares there, because consecutive squares near $x$ are about $2\sqrt x$ apart, and at most one $k$-th power for each $k\ge3$, because consecutive $k$-th powers are farther than $h$ apart for large $x$. They contribute $O((\log^3x+\log x)\log x)=o(h)$. So $\theta(x+h)-\theta(x)>0$: there is a prime.
\end{solution}

\begin{solution}{ex:vonkoch}
\idea The passage from $\psi$ to $\pi$ divides the main term by $\log x$. The error that comes out is $O(\sqrt x\log x)$, one logarithm smaller than the error for $\psi$.
\writeup Under the Riemann hypothesis, Theorem~\ref{thm:rhequiv} and Chapter~\ref{ch:mertens} give
\[
\theta(x)-x=\psi(x)-x+O(\sqrt x\log x)=O(\sqrt x\log^2x).
\]
In the identity of the proof of Corollary~\ref{cor:pnterrorpi}, the first term is $O(\sqrt x\log x)$, the constant is $O(1)$, and the integral is
\[
O\Bigl(\int_2^x\frac{\sqrt t\log^2t}{t\log^2t}\dd t\Bigr)=O(\sqrt x).
\]
So $\pi(x)-\operatorname{Li}(x)=O(\sqrt x\log x)$. The bound for $\psi$ in Theorem~\ref{thm:rhequiv} remains $O(\sqrt x\log^2x)$.
\end{solution}

\begin{solution}{ex:bias}
\idea Prime powers of exponent $2$ all fall into the square classes. The identity moves those squares onto one side as a term $\theta(\sqrt x)$. The sums $\psi(x,\chi)$ stay in the formula, and this exercise does not estimate them.
\writeup (a) One has $\chi_{-3}(p^2)=1$ for $p\ne3$ and $\chi_{-3}(9)=0$, so the squares contribute $\theta(\sqrt x)-\log3$ to $\psi(x,\chi_{-3})$ when $x\ge9$. The primes contribute $\theta(x;3,1)-\theta(x;3,2)$, and the higher powers contribute $O(x^{1/3}\log x)$. Rearranging gives the identity. The term $\theta(\sqrt x)$ is a head start for the class $2$, the class that does not contain the squares.

(b) Let $\theta(x,\chi)=\sum_{p\le x}\chi(p)\log p$, so $\theta(x,\chi_0)=\theta(x)-\log2$. By orthogonality (Chapter~\ref{ch:characters}), $\theta(x;8,a)=\frac14\sum_\chi\chi(a)\theta(x,\chi)$, and the characters modulo $8$ are real. For $\chi\ne\chi_0$, the same grouping as in Proposition~\ref{prop:biasidentity} gives
\[
\psi(x,\chi)=\theta(x,\chi)+\theta(\sqrt x)-\log2+O(x^{1/3}\log x),
\]
since $\chi(p^2)=1$ for odd $p$. Hence
\[
\begin{aligned}\theta(x;8,a)&=\frac{\theta(x)}4+\frac14\sum_{\chi\ne\chi_0}\chi(a)\psi(x,\chi)\eqbr-\frac{\theta(\sqrt x)}4\sum_{\chi\ne\chi_0}\chi(a)+O(x^{1/3}\log x).\end{aligned}
\]
By orthogonality, $\sum_\chi\chi(a)=4\cdot\mathbf1_{a\equiv1}$, so $\sum_{\chi\ne\chi_0}\chi(a)=3$ if $a\equiv1$ and equals $-1$ otherwise. The coefficient of $-\theta(\sqrt x)/4$ is therefore $3$ for $a=1$ and $-1$ for $a\ne1$. Class $1$ carries the term $-\frac34\theta(\sqrt x)$. Each of the other three classes carries $+\frac14\theta(\sqrt x)$. Since $\theta(\sqrt x)\sim\sqrt x$, those four coefficients are $-\frac34,\frac14,\frac14,\frac14$. The sums $\psi(x,\chi)$ are still present. Nothing in the identity bounds them, so the identity alone does not assign a sign to $\theta(x;8,a)-\theta(x)/4$ at a given $x$.

At $x=10^6$ the four numbers $(\theta(x;8,a)-\theta(x)/4)/\sqrt x$ are $-0.838$, $0.336$, $-0.071$ and $0.572$ for $a=1,3,5,7$. Averaged with respect to $\mathrm dx/x$ over $10^3\le x\le10^7$, they are $-0.719$, $0.234$, $0.247$ and $0.233$. The handicap coefficients are $-0.75$, $0.25$, $0.25$ and $0.25$. At the single point $x=10^6$ the four numbers sit away from those coefficients, so the character sums are visible in the difference. The average lies nearer the coefficients. That comparison is a numerical observation about the whole left side. It is not an estimate of $\psi(x,\chi)$.
\end{solution}

\begin{solution}{ex:cesaro}
\idea The weight $1-n/x$, in place of the sharp cut-off, replaces $1/s$ by $1/(s(s+1))$ in Perron's formula. The extra decay in $|t|$ makes the integral absolutely convergent, so no truncation error appears, and the residue at each zero gains a factor $1/(\rho+1)$.
\writeup \emph{The kernel.} For $y>1$, shift $\int_{c-iT}^{c+iT} y^s/(s(s+1))\dd s$ to the left as in Lemma~\ref{lem:perron}. The residues at $0$ and $-1$ are $1$ and $-1/y$, and the other sides are $O(y^c/T^2)$, which tend to $0$. For $y<1$, shift right: there is no residue. For $y=1$, both sides are continuous in $y$ (the integrand is bounded by $\max(y,1)^c/|s|^2$). So the kernel is $(1-1/y)\mathbf1_{y\ge1}$. With $y=x/n$ and term-by-term integration on $\operatorname{Re}s=c>1$,
\[
\begin{aligned}\sum_{n\le x}\Lambda(n)\Bigl(1-\frac nx\Bigr)&=\frac1{2\pi i}\int_{c-i\infty}^{c+i\infty}\Bigl(-\frac{\zeta'}\zeta(s)\Bigr)\eqbr\quad\times\frac{x^s}{s(s+1)}\dd s.\end{aligned}
\]
\emph{The shift.} Take the rectangle of Theorem~\ref{thm:explicit} with a good height $T'$. The residues are: $x/2$ at $1$; $-x^\rho/(\rho(\rho+1))$ at each zero; $-\zeta'/\zeta(0)=-\log2\pi$ at $0$; and
\[
-\frac{\zeta'}\zeta(-1)\cdot\frac{x^{-1}}{-1}=\frac{\zeta'(-1)}{\zeta(-1)\,x}
\]
at $-1$, where $\zeta(-1)=-\frac1{12}\ne0$. At $-2k$ for $k\ge1$ the residue is
\[
-\frac{x^{-2k}}{(-2k)(-2k+1)}=-\frac{x^{-2k}}{2k(2k-1)}.
\]
The left side tends to $0$ as before. The horizontal sides are now $O(x\log^2T/(T^2\log x))$ by Lemmas~\ref{lem:goodT} and~\ref{lem:leftzeta}, and the vertical tail $|t|>T'$ is $O(x\log x/T)$, since $-\zeta'/\zeta(c+it)$ is bounded by $\log x+A$. Everything tends to $0$ as $T\to\infty$. Finally
\[
\sum_\rho\frac1{|\rho(\rho+1)|}\le\sum_\rho\frac1{|\rho|^2}<\infty
\]
(Exercise~\ref{ex:zerosums} gives much more), and $|x^\rho|\le x$, so the series converges absolutely.
\end{solution}

\begin{solution}{ex:halfjump}
\idea At a prime the definition of $\psi_0$ removes half the jump. Corollary~\ref{cor:exactexplicit} then solves for the sum over zeros, because every other term is continuous and explicit.
\writeup The prime powers up to $5$ are $2,3,4,5$, so
\[
\psi(5)=\log2+\log3+\log2+\log5=\log60=4.0943.
\]
Since $5$ is prime, $\Lambda(5)=\log5$ and
\[
\psi_0(5)=\psi(5)-\frac12\log5=\log60-\frac12\log5=3.2896.
\]
The difference is $\frac12\log5=0.8047$. Now $1-1/25=24/25=0.96$ and $\frac12\log0.96=-0.02041$, so
\[
5-\log(2\pi)-\frac12\log\Bigl(1-\frac1{25}\Bigr)=5-1.8379+0.02041=3.1825.
\]
Corollary~\ref{cor:exactexplicit} says this quantity minus the symmetric sum equals $\psi_0(5)$. The symmetric sum therefore equals $3.1825-3.2896=-0.1071$.
\end{solution}

\begin{solution}{ex:trivten}
\idea The trivial zeros are a geometric series in the variable $x^{-2}$. At $x=10$ that series is a few thousandths, beside a main term of size $10$.
\writeup For $|z|<1$, one has $\sum_{k\ge1}z^k/k=-\log(1-z)$. Set $z=x^{-2}$. Then
\[
\sum_{k\ge1}\frac{x^{-2k}}{2k}=\frac12\sum_{k\ge1}\frac{(x^{-2})^k}k=-\frac12\log(1-x^{-2}).
\]
At $x=10$, the sum equals $-\frac12\log0.99=0.00503$, to three significant figures. In Corollary~\ref{cor:exactexplicit} this contribution sits beside the main term $x=10$.
\end{solution}

\begin{solution}{ex:balanceT}
\idea The zero-free region contributes the denominator $\log(T+2)$. The choice $\log T=\sqrt{\log x}$ makes $\log T$ and $\log(T+2)$ almost equal, so replacing the ratio by $\frac12$ is legal for large $x$, and at $\log x=100$ the ratio is already indistinguishable from $1$ at three decimal places.
\writeup (a) With $\log T=\sqrt{\log x}$ one has $T\to\infty$ as $x\to\infty$, and
\[
\frac{\log T}{\log(T+2)}=\frac{\log T}{\log T+\log(1+2/T)}.
\]
Since $\log(1+2/T)\sim2/T\to0$, the ratio tends to $1$.

(b) If $\log x=100$, then $\log T=10$ and $T=e^{10}=22026.4658$. Then $T+2=22028.4658$ and $\log(T+2)=10.000091$, so
\[
\frac{\log T}{\log(T+2)}=0.999991,
\]
which is greater than $\frac12$. To three decimal places the ratio rounds to $1.000$.

(c) One has $\log x/\log(T+2)=100/10.000091=9.99991$. That is $\sqrt{\log x}$ times the ratio in (b): $\sqrt{100}=10$ and $10\times0.999991=9.99991$. The factor $\frac12$ used in the proof of Theorem~\ref{thm:pnterror} replaces $9.99991$ by $5$. It is legal, and at this $x$ it is generous.
\end{solution}

\begin{solution}{ex:psi0koch}
\idea The two functions differ by at most half a logarithm. Every error term in Theorem~\ref{thm:rhequiv} is much larger than $\log x$, so the equivalence is unchanged, and the power of $\log x$ in (b) stays $2$.
\writeup If $x$ is not an integer, then $\psi_0(x)=\psi(x)$. If $x$ is an integer, then $\psi(x)-\psi_0(x)=\frac12\Lambda(x)$. And $\Lambda(x)\le\log x$, so $|\psi(x)-\psi_0(x)|\le\frac12\log x$ for $x\ge2$.

Assume (a), the Riemann hypothesis. The proof of (a)$\Rightarrow$(b) in Theorem~\ref{thm:rhequiv} bounds $\psi_0(x)-x$ by $O(\sqrt x\log^2x)$, because the sum of $|x^\rho/\rho|$ is $O(\sqrt x\log^2x)$ and the truncated error is $O(\log^2x)$. Adding the difference $\psi-\psi_0$, which is $O(\log x)$, does not change that bound. The same proof with the roles reversed, starting from a bound for $\psi_0$, absorbs $O(\log x)$ into $O(\sqrt x\log^2x)$. So (a) is equivalent to $\psi_0(x)=x+O(\sqrt x\log^2x)$.

For (c), the difference $O(\log x)$ is absorbed in $O(x^{1/2+\varepsilon})$ as well. Replacing $\psi$ by $\psi_0$ does not improve the power of $\log x$ in (b): the sum $\sum_{|\gamma|<x}1/|\rho|$ in Lemma~\ref{lem:sumrho} is what produces $\log^2x$, and removing half a jump removes only $\log x$.
\end{solution}

\section*{Sources and further reading}

\begin{itemize}
\item H.~von Mangoldt, \emph{Zu Riemanns Abhandlung ``\"Uber die Anzahl der Primzahlen unter einer gegebenen Gr\"osse''}, Journal f\"ur die reine und angewandte Mathematik 114 (1895), 255--305. The proof of the explicit formula used in this chapter.
\item C.-J.~de la Vall\'ee Poussin, \emph{Sur la fonction $\zeta(s)$ de Riemann et le nombre des nombres premiers inf\'erieurs \`a une limite donn\'ee}, M\'emoires couronn\'es de l'Acad\'emie royale de Belgique 59 (1899).
\item H.~von Koch, \emph{Sur la distribution des nombres premiers}, Acta Mathematica 24 (1901), 159--182.
\item J.~E.~Littlewood, \emph{Sur la distribution des nombres premiers}, Comptes Rendus de l'Acad\'emie des Sciences 158 (1914), 1869--1872.
\item H.~Cram\'er, \emph{Some theorems concerning prime numbers}, Arkiv f\"or Matematik, Astronomi och Fysik 15 (1920), 1--33.
\item J.~Leech, \emph{Note on the distribution of prime numbers}, Journal of the London Mathematical Society 32 (1957), 56--58.
\item L.~Schoenfeld, \emph{Sharper bounds for the Chebyshev functions $\theta(x)$ and $\psi(x)$. II}, Mathematics of Computation 30 (1976), 337--360.
\item M.~Rubinstein and P.~Sarnak, \emph{Chebyshev's bias}, Experimental Mathematics 3 (1994), 173--197.
\item A.~Granville and G.~Martin, \emph{Prime number races}, American Mathematical Monthly 113 (2006), 1--33. A survey of the races, and a good next thing to read on Chebyshev's bias.
\item H.~Davenport, \emph{Multiplicative Number Theory}, 3rd ed., Springer GTM~74, 2000, Chapters~17--20 (the explicit formulae for $\zeta$ and $L(s,\chi)$, the prime number theorem with an error term).
\item H.~L.~Montgomery and R.~C.~Vaughan, \emph{Multiplicative Number Theory I}, Cambridge, 2007, Chapters~12, 13 and~15 (explicit formulae, conditional estimates, oscillation theorems).
\item H.~M.~Edwards, \emph{Riemann's Zeta Function}, Academic Press, 1974, Chapters~1 and~3 (Riemann's formula for $\pi(x)$).
\item The epigraph is a sentence widely repeated as Hilbert's. No primary source for it is known, which is why the line is marked ``attributed''.
\end{itemize}
